\section{The Vorton-Scalar and the deflation clock}
\label{sec:clock}

\subsection{The model}

The model contains, besides the fermion field, a complex scalar field, the \emph{Vorton-Scalar}
$\PhiVS$, whose vacuum value gives the fermions a mass and whose exchange binds two fermions into a
pair. This subsection states the complete model on the deflating spacetime: the fermion field with a
mass term and with the self-interaction of Sec.~\ref{sec:dirac}, and the Vorton-Scalar. The
Vorton-Scalar brings three real, positive constants: its self-coupling $\lamVS$, the constant
$\vVS$ that fixes the magnitude $\vVS/\sqrt2$ at which the potential of the scalar is lowest, and
its Yukawa coupling $\yVS$ to the fermions. The fermion field brings the mass $m_D$ of its
Lagrangian, called the \emph{bare Dirac mass} from here on, and the couplings $g_r$ of its
self-interaction. The chirality matrix is denoted by $\chi$, with
\SIG{$\chi=\gamma^{0}\gamma^{1}\cdots\gamma^{7}$, which is real on the $4+4$ spacetime,}{$\chi=i\,\gamma^{0}\gamma^{1}\cdots\gamma^{7}$,}
and the chiral projectors are $P^{L,R}=\tfrac12(1\mp\chi)$.

On a curved spacetime a Lagrangian is multiplied by the volume factor $\Vd$, derivatives of spinors
become covariant derivatives, and the flat metric is replaced by the curved one.
%%<1671>
The Lagrangian of the model is
\begin{align}
\mathcal L_{\tg}^{\mathrm{VS}}=\Vd\,\Big[&\bar\psi\,\big(i\gamma^{\mu}\nabla_\mu-m_D\big)\psi+\Lself\nonumber\\
&-\gmup{\mu\nu}\,\partial_\mu\PhiVS^{*}\,\partial_\nu\PhiVS\nonumber\\
&-\lamVS\Big(|\PhiVS|^{2}-\frac{\vVS^{2}}{2}\Big)^{2}\nonumber\\
&-\yVS\,\bar\psi\big(\PhiVS\,P^{R}+\PhiVS^{*}\,P^{L}\big)\psi\Big],
\label{eq:LVS}
\end{align}
with $\bar\psi=\psi^\dagger\beta$, $\beta=\gamma^{4}$. The first line is the Lagrangian of
Sec.~\ref{sec:dirac} with the bare Dirac mass; the second is the kinetic term of the scalar (the minus
sign belongs to the sign convention of the metric, in which the component of the time is negative,
so that the term contains $+|\partial_4\PhiVS|^{2}$); the third is its potential; the fourth is the
Yukawa coupling.
%%</1671>
%%<1644>
On the $4+4$ spacetime the Yukawa coupling can be written only in the sector of fields that do not
depend on the additional times, the quantized sector of Sec.~\ref{sec:forms}, because only there
does the adjoint $\bar\psi=\psi^\dagger\beta$ exist with which the chiral projectors combine into
a real, invariant term; the model is defined there. On the deflating spacetime Sec.~\ref{sec:modes}
gives that sector a dynamical reason. In it the model is
\begin{align}
\mathcal L_{\tg}^{\mathrm{VS}}={}&\hpsi^\dagger\big(i\partial_4+i\sum_{j=0}^{3}\alpha^{j}\partial_j-\beta\,m_D\big)\hpsi+\Vd\,\Lself\nonumber\\
&+\Vd\,\Big[|\partial_4\PhiVS|^{2}-\sum_{j=0}^{3}|\partial_j\PhiVS|^{2}\Big]\nonumber\\
&-\Vd\,\lamVS\Big(|\PhiVS|^{2}-\frac{\vVS^{2}}{2}\Big)^{2}\nonumber\\
&-\Vd\,\yVS\,\bar\psi\big(\PhiVS\,P^{R}+\PhiVS^{*}\,P^{L}\big)\psi ,
\label{eq:LVS}
\end{align}
with the Hermitian matrices $\beta=-i\gamma^{4}$ and $\alpha^{j}=-\gamma^{4}\gamma^{j}$,
$\bar\psi=\psi^\dagger\beta$, and the self-interaction $\Lself$ in the quantized-sector form
(\ref{eq:LselfQ}). The first line is the positive Lagrangian (\ref{eq:Lpos}) for the weighted field
$\hpsi=\sqrt{\Vd}\,\psi$, with the bare Dirac mass, and the self-interaction; without the
self-interaction its equation of motion is the covariant Dirac equation (\ref{eq:dirac}) in this
sector. The other terms carry the volume factor.
%%</1644>
The Lagrangian is invariant under the phase transformation $\psi\to e^{i\epsilon}\psi$ of the
fermion number, $U(1)_F$, whatever its constants are. Without the bare Dirac mass, and with a
self-interaction that satisfies the condition (\ref{eq:chiralcond}), it is also invariant under the
chiral transformation $\UVS$: $\psi\to e^{i\epsilon\chi}\psi$ with
$\PhiVS\to e^{-2i\epsilon}\PhiVS$. These statements hold on the deflating spacetime as on flat
spacetime, because the transformations do not involve the coordinates. \emph{The case $m_D=0$,
$g_r=0$ is the model in which the whole mass of the fermions comes from the Vorton-Scalar and the
whole interaction from its exchange;} Secs.~\ref{sec:gap} to \ref{sec:order} treat it first, and
Secs.~\ref{sec:ks}, \ref{sec:gapcontact} and \ref{sec:cjtcontact} then add the bare mass and the
self-interaction and say what each changes.

A constant field $\PhiVS=\vVS/\sqrt2$ with $\psi=0$ solves the field equations on the deflating
spacetime, because a constant has no derivatives: the vacuum value is the same at all times. It
breaks $\UVS$. Around it the Vorton-Scalar has two modes: the \emph{radial mode}, a real field of the
mass $m_s$, and the \emph{phase mode}, a real massless field, the Goldstone boson of the broken
symmetry (massless if $\UVS$ is exact; Sec.~\ref{sec:cjtcontact} computes the mass that a bare Dirac
mass gives it). The fermions acquire the mass $m=m_D+m_Y$, of which $m_Y=\yVS\vVS/\sqrt2$ comes from
the Vorton-Scalar. Both modes couple to the fermions with the same strength
$\gVS$, the radial mode to the scalar density $\bar\psi\psi$ and the phase mode to the pseudoscalar
density $\bar\psi\,i\chi\,\psi$:
\begin{equation}
\begin{aligned}
m_Y&=\frac{\yVS\,\vVS}{\sqrt2},&\qquad m_s&=\sqrt{2\lamVS}\,\vVS,\\
\gVS&=\frac{m_Y}{\vVS},&\qquad\aVS&=\frac{\gVS^{2}}{4\pi}.
\end{aligned}
\label{eq:masses}
\end{equation}

\subsection{Fields that do not depend on the extra coordinates}

Consider fields that do not depend on the four extra coordinates
$x0,x5,x6,x7$, and introduce the weighted fields
\begin{equation}
\hpsi=\sqrt{\Vd}\,\psi,\qquad\hat\Phi_{\mathrm{VS}}=\sqrt{\Vd}\,\PhiVS .
\end{equation}
Each term of Eq.~(\ref{eq:LVS}) is rewritten in turn. The kinetic term and the mass term of the
fermions are the flat ones for $\hpsi$ (Sec.~\ref{sec:dirac}). The self-interaction contains four
fermion fields and therefore the factor $\Vd\cdot\Vd^{-2}$, as in Eq.~(\ref{eq:grun}). The Yukawa
term contains one scalar and two fermion fields
and therefore the factor $\Vd\cdot\Vd^{-3/2}$. The potential becomes
$(\lamVS/\Vd)\,(|\hat\Phi_{\mathrm{VS}}|^{2}-\Vd\vVS^{2}/2)^{2}$. In the kinetic term of the scalar the
derivative along the time acts on the weight as well,
$\sqrt{\Vd}\,\partial_4\PhiVS=\partial_4\hat\Phi_{\mathrm{VS}}+\tfrac12\gsum\,\hat\Phi_{\mathrm{VS}}$,
which gives a term $\tfrac14\gsum^{2}|\hat\Phi_{\mathrm{VS}}|^{2}$ and a total derivative. The result
is the flat Lagrangian for the weighted fields, with four replacements and one additional term:
\begin{equation}
\begin{aligned}
&\yVS\to\frac{\yVS}{\sqrt{\Vd(x4)}},\qquad\lamVS\to\frac{\lamVS}{\Vd(x4)},\qquad g_r\to\frac{g_r}{\Vd(x4)},\\
&\vVS\to\vVS\sqrt{\Vd(x4)},\qquad\text{and }+\frac{\gsum^{2}}{4}\,|\hat\Phi_{\mathrm{VS}}|^{2}.
\end{aligned}
\label{eq:couplings}
\end{equation}
The bare Dirac mass $m_D$ is not replaced. With a constant volume of the extra coordinates the same
replacements define the couplings of a four-dimensional description on flat spacetime. Here that
volume is the coordinate volume of a cell times $\Vd(x4)$, and it shrinks. The additional term lowers the square of
the mass of each mode of the scalar by $\gsum^{2}/4$ in the equation of the weighted field;
Section~\ref{sec:scalarmodes} shows what it does. From here on the couplings $\yVS$, $\lamVS$,
$\vVS$, $\gVS$, $\aVS$ and $g_r$ are those of the four-dimensional description at the time
considered.

\subsection{The deflation clock}

Define the \emph{deflation clock}
\begin{equation}
\td=\gsum\,(x4-x4_{\mathrm{ref}}),
\label{eq:clock}
\end{equation}
the number of e-foldings by which the volume $\Vd$ has shrunk since a reference time
$x4_{\mathrm{ref}}$, which is chosen in each application. By Eq.~(\ref{eq:couplings}) every quantity
of the four-dimensional description is proportional to $e^{n\td}$ with a definite exponent $n$
(Table~\ref{tab:exponents}). Three observations organize what follows.

\begin{table*}
\caption{\label{tab:exponents}How the quantities of the four-dimensional description follow the
deflation clock: each is proportional to $e^{n\td}$. The last two rows hold under additional
assumptions, which Secs.~\ref{sec:cosmo} and \ref{sec:higgs} state and examine.}
\begin{ruledtabular}
\begin{tabular}{p{0.37\textwidth}p{0.47\textwidth}c}
quantity & meaning & exponent $n$ \\
\hline
{{ROWS:tbl_exponents}}
\end{tabular}
\end{ruledtabular}
\end{table*}

\emph{The masses do not change.} The bare Dirac mass is not affected at all, and the part $m_Y$ of
the fermion mass and the mass $m_s$ of the radial mode contain the couplings in combinations from
which the volume cancels. So do the Fermi momentum and the Fermi
energy of matter with a given number of particles per unit of ordinary volume.

\emph{The couplings grow.} The coupling $\gVS=m_Y/\vVS$ grows as $e^{\td/2}$, because the
four-dimensional vacuum value falls. The strength of the attraction, $\aVS$, grows as $e^{\td}$: it
doubles every time the volume of the deflating coordinates halves. Every coupling $g_r$ of the
self-interaction, and with them the couplings $\GS$, $\GV$ of Sec.~\ref{sec:ks}, grows as $e^{\td}$
as well: \emph{the contact interaction and the exchange of the Vorton-Scalar keep their ratio.}

\emph{The two numbers that decide the state of matter grow at the same rate.} These are
\begin{equation}
\DVS=\aVS\,\frac{m}{m_s},\qquad C=\frac{\gVS^{2}m^{2}}{m_s^{2}} .
\label{eq:DC}
\end{equation}
The first decides whether two fermions form a bound pair (Sec.~\ref{sec:pair}), the second whether
matter keeps its fermion mass (Sec.~\ref{sec:cjt}). Their ratio $C/\DVS=4\pi\,m/m_s$ is fixed.

On flat spacetime the model is specified by the masses and the couplings. On a deflating spacetime
it is specified by the masses and the values of the couplings \emph{at one time}; the deflation then
carries it through all larger couplings. The constants of the flat theory become functions of the
time, and its thresholds are passed one after another (Fig.~\ref{fig:clock}).

\begin{figure*}
\includegraphics[width=0.9\textwidth]{fig_clock}
\caption{\label{fig:clock}The deflation clock. (a) How the quantities of the four-dimensional
description change with the clock $\td$ (logarithmic scale; values relative to $\td=0$). (b) The
order of events along the deflation for six ratios $m_s/m$: the clock reading of each event relative
to the moment at which two fermions form a bound pair (Table~\ref{tab:events}).}
\end{figure*}

\subsection{Modes of the Vorton-Scalar}
\label{sec:scalarmodes}

Let $\varphi_{\mathrm{bos}}$ stand for the radial mode or the phase mode, with the mass
$m_{\mathrm{bos}}=m_s$ or $m_{\mathrm{bos}}=0$. To first order in the field, a mode with a wavenumber
of the magnitude $k$ along the flat coordinates and the wavenumbers $k_t$ along the deflating ones
obeys
\begin{equation}
\begin{aligned}
&\frac{1}{\Vd}\,\partial_4\big(\Vd\,\partial_4\varphi_{\mathrm{bos}}\big)\\
&\quad+\Big(k^{2}\SIG{-}{+}\sum_t\frac{k_t^{2}}{b_t^{2}}+m_{\mathrm{bos}}^{2}\Big)\varphi_{\mathrm{bos}}=0 .
\end{aligned}
\label{eq:scalar}
\end{equation}
\SIG{The wavenumbers along the deflating coordinates, which are time-like, enter with the sign
opposite to that of $k^{2}$, as for the fermions.}{}

\emph{Modes that do not depend on the deflating coordinates.} With $k_t=0$ the equation has constant
coefficients, and its solutions are the exponentials
\begin{equation}
\varphi_{\mathrm{bos}}\propto\exp\Big\{\Big[\frac{\gsum}{2}\pm\sqrt{\frac{\gsum^{2}}{4}-k^{2}-m_{\mathrm{bos}}^{2}}\,\Big]x4\Big\}.
\label{eq:exponents}
\end{equation}
If $k^{2}+m_{\mathrm{bos}}^{2}>\gsum^{2}/4$, the mode oscillates with the frequency
$\varpi_b=(k^{2}+m_{\mathrm{bos}}^{2}-\gsum^{2}/4)^{1/2}$, and its amplitude grows as
$e^{\gsum\,x4/2}$. The growth is that of the weight alone: the weighted field has a constant
amplitude, and the number of quanta in a cell of fixed coordinates is conserved while the cell
shrinks. Otherwise the mode does not oscillate. For the radial mode this never happens if
$m_s>\gsum/2$, which holds by many orders of magnitude for any rate that observation allows
(Sec.~\ref{sec:cosmo}). For the phase mode it happens at wavelengths longer than $4\pi/\gsum$. At
$k=0$ the two exponents of the phase mode are $0$ and $\gsum$: a constant phase, which costs nothing,
and a phase that advances at a rate proportional to $1/\Vd$, which is the conservation of the charge
of $\UVS$ in a shrinking volume. Neither is an instability of the vacuum. Figure~\ref{fig:scalar}
shows the exponents.

\begin{figure*}
\includegraphics[width=\textwidth]{fig_scalar}
\caption{\label{fig:scalar}Modes of the Vorton-Scalar that do not depend on the deflating
coordinates. (a) Real parts of the two exponents of Eq.~(\ref{eq:exponents}) against the wavenumber,
for the phase mode and for a radial mode, in units of $\gsum$. (b) Their imaginary part, the
frequency of oscillation. (c) The field that a static source builds up from zero, divided by the form
$e^{\td}/(k^{2}+m_{\mathrm{bos}}^{2})$ that leads to Eq.~(\ref{eq:yukawa}).}
\end{figure*}

\emph{The potential between two fermions.} Two fermions at rest attract each other by exchanging the
radial mode. The density of a fixed number of fermions grows as $1/\Vd$, so that the radial mode
obeys Eq.~(\ref{eq:scalar}) with a source proportional to $e^{\gsum\,x4}$. The solution that follows
the source is $e^{\gsum\,x4}/(k^{2}+m_{\mathrm{bos}}^{2})$ times the source, as substitution shows:
the terms $\gsum^{2}-\gsum\cdot\gsum$ cancel, and the denominator contains the \emph{unshifted} mass.
The potential between two fermions at rest is therefore, at every moment, a Yukawa potential,
\begin{equation}
V_s(r)=-\aVS(x4)\,\frac{e^{-m_sr}}{r},
\label{eq:yukawa}
\end{equation}
with the range $1/m_s$ unchanged and the strength of the moment. The numerical solution from rest
approaches this form with a transient that decays as $e^{-\td/2}$ [Fig.~\ref{fig:scalar}(c)]; after
$\td=20$ the relative difference is below $4\times10^{-5}$ in the four cases computed.

\emph{Modes that depend on a deflating coordinate.} For $k_t\ne0$ the weighted field obeys
$\partial_4^{2}\hat\varphi_{\mathrm{bos}}+(\varpi_b^{2}\SIG{-}{+}\kd^{2})\,\hat\varphi_{\mathrm{bos}}=0$
with the physical wavenumber $\kd$ of Eq.~(\ref{eq:kd}), and the scalar behaves as the fermion field
does.
%%<1671>
The frequency grows without bound. When one rate governs, the solutions are Bessel functions of the
imaginary order $i\varpi_b/\rt$ and of the argument $\ud$, and a mode that starts empty ends with
\begin{equation}
\ncr=\frac{1}{e^{2\pi\varpi_b/\rt}-1}
\label{eq:bose}
\end{equation}
quanta: the Bose-Einstein distribution at the same temperature $\Tdefl$ as for the fermions. The
program confirms Eq.~(\ref{eq:bose}) in 36 cases to a relative deviation of $1.9\times10^{-5}$
(Fig.~\ref{fig:creation}).
%%</1671>
%%<1644>
Beyond the turning value, at which $\kd=\varpi_b$, one solution decays and one grows. The bounded
solution is the modified Bessel function of the imaginary order $i\varpi_b/\rt$ that decays at large
argument. At early times it is a standing wave; at late times it falls as $e^{-\ud}$
(Fig.~\ref{fig:extinction}). The scalar field, like the fermion field, loses its dependence on the
deflating times.
%%</1644>

